% Part: second-order-logic
% Chapter: sol-and-set-theory
% Section: comparing-sets

\documentclass[../../../include/open-logic-section]{subfiles}

\begin{document}

\olfileid{sol}{set}{cmp}

\olsection{مقارنة المجموعات}

\begin{prop}
!!{formula}~$\lforall[{\OLId{latin-ordinary}{x}}][({\OLId{latin-looped}{X}}({\OLId{latin-ordinary}{x}}) \lif {\OLId{latin-looped}{Y}}({\OLId{latin-ordinary}{x}}))]$ تعريف لعلاقة الاحتواء
الجزئي؛ فـ$\Sat{{\OLId{latin-looped}{M}}}{\lforall[{\OLId{latin-ordinary}{x}}][({\OLId{latin-looped}{X}}({\OLId{latin-ordinary}{x}}) \lif {\OLId{latin-looped}{Y}}({\OLId{latin-ordinary}{x}}))]}[{\OLId{latin-ordinary}{s}}]$ إذا وفقط
إذا~${\OLId{latin-ordinary}{s}}({\OLId{latin-looped}{X}}) \subseteq {\OLId{latin-ordinary}{s}}({\OLId{latin-looped}{Y}})$.
\end{prop}

\begin{prop}
!!{formula}~$\lforall[{\OLId{latin-ordinary}{x}}][({\OLId{latin-looped}{X}}({\OLId{latin-ordinary}{x}}) \liff {\OLId{latin-looped}{Y}}({\OLId{latin-ordinary}{x}}))]$ تعريف لهوية المجموعات؛
فـ$\Sat{{\OLId{latin-looped}{M}}}{\lforall[{\OLId{latin-ordinary}{x}}][({\OLId{latin-looped}{X}}({\OLId{latin-ordinary}{x}}) \liff {\OLId{latin-looped}{Y}}({\OLId{latin-ordinary}{x}}))]}[{\OLId{latin-ordinary}{s}}]$ إذا وفقط
إذا~${\OLId{latin-ordinary}{s}}({\OLId{latin-looped}{X}}) = {\OLId{latin-ordinary}{s}}({\OLId{latin-looped}{Y}})$.
\end{prop}

\begin{prop}
!!{formula}~$\lexists[{\OLId{latin-ordinary}{x}}][{\OLId{latin-looped}{X}}({\OLId{latin-ordinary}{x}})]$ تعريف لعدم الخلو؛
فـ$\Sat{{\OLId{latin-looped}{M}}}{\lexists[{\OLId{latin-ordinary}{x}}][{\OLId{latin-looped}{X}}({\OLId{latin-ordinary}{x}})]}[{\OLId{latin-ordinary}{s}}]$ إذا وفقط إذا~${\OLId{latin-ordinary}{s}}({\OLId{latin-looped}{X}}) \neq \emptyset$.
\end{prop}

قولنا إن المجموعة~${\OLId{latin-looped}{X}}$ لا تزيد حجمًا على المجموعة~${\OLId{latin-looped}{Y}}$، أي
$\cardle{{\OLId{latin-looped}{X}}}{{\OLId{latin-looped}{Y}}}$، معناه وجود دالة !!a{injective}~${\OLId{latin-ordinary}{f}}\colon {\OLId{latin-looped}{X}} \to {\OLId{latin-looped}{Y}}$.
ولما أمكن التعبير عن تباين الدالة وعن وقوع قيمها عند عناصر~${\OLId{latin-looped}{X}}$
في~${\OLId{latin-looped}{Y}}$، أمكن تعريف علاقة عدم الزيادة في الحجم بين مجموعات المجال
الجزئية.

\begin{prop}
الصيغة
\[
\lexists[{\OLId{latin-ordinary}{u}}][(\lforall[{\OLId{latin-ordinary}{x}}][({\OLId{latin-looped}{X}}({\OLId{latin-ordinary}{x}}) \lif {\OLId{latin-looped}{Y}}({\OLId{latin-ordinary}{u}}({\OLId{latin-ordinary}{x}})))] \land \lforall[{\OLId{latin-ordinary}{x}}][\lforall[{\OLId{latin-ordinary}{y}}][(\eq[{\OLId{latin-ordinary}{u}}({\OLId{latin-ordinary}{x}})][{\OLId{latin-ordinary}{u}}({\OLId{latin-ordinary}{y}})] \lif
      \eq[{\OLId{latin-ordinary}{x}}][{\OLId{latin-ordinary}{y}}])]])]
\]
تعريف لعلاقة عدم الزيادة في الحجم.
\end{prop}

ومجموعتان متساويتان حجمًا، أو «متكافئتان عدديًا»،
$\cardeq{{\OLId{latin-looped}{X}}}{{\OLId{latin-looped}{Y}}}$، إذا وفقط إذا وجدت بينهما دالة
!!a{bijective}~${\OLId{latin-ordinary}{f}}\colon {\OLId{latin-looped}{X}} \to {\OLId{latin-looped}{Y}}$.

\begin{prop}
الصيغة
\begin{multline*}
  \lexists[{\OLId{latin-ordinary}{u}}][(\lforall[{\OLId{latin-ordinary}{x}}][({\OLId{latin-looped}{X}}({\OLId{latin-ordinary}{x}}) \lif {\OLId{latin-looped}{Y}}({\OLId{latin-ordinary}{u}}({\OLId{latin-ordinary}{x}})))] \land {}\\
    \lforall[{\OLId{latin-ordinary}{x}}][\lforall[{\OLId{latin-ordinary}{y}}][(({\OLId{latin-looped}{X}}({\OLId{latin-ordinary}{x}}) \land {\OLId{latin-looped}{X}}({\OLId{latin-ordinary}{y}}) \land \eq[{\OLId{latin-ordinary}{u}}({\OLId{latin-ordinary}{x}})][{\OLId{latin-ordinary}{u}}({\OLId{latin-ordinary}{y}})]) \lif \eq[{\OLId{latin-ordinary}{x}}][{\OLId{latin-ordinary}{y}}])]]
  \land {} \\\lforall[{\OLId{latin-ordinary}{y}}][({\OLId{latin-looped}{Y}}({\OLId{latin-ordinary}{y}}) \lif \lexists[{\OLId{latin-ordinary}{x}}][({\OLId{latin-looped}{X}}({\OLId{latin-ordinary}{x}})
      \land \eq[{\OLId{latin-ordinary}{y}}][{\OLId{latin-ordinary}{u}}({\OLId{latin-ordinary}{x}})])])])]
\end{multline*}
تعريف لعلاقة التكافؤ العددي.
\end{prop}

وسنستعمل اختصارًا لهذه !!{formula}s الرموز
${\OLId{latin-looped}{X}} \subseteq {\OLId{latin-looped}{Y}}$، و${\OLId{latin-looped}{X}} = {\OLId{latin-looped}{Y}}$، و${\OLId{latin-looped}{X}} \neq \emptyset$، و$\cardle{{\OLId{latin-looped}{X}}}{{\OLId{latin-looped}{Y}}}$،
و$\cardeq{{\OLId{latin-looped}{X}}}{{\OLId{latin-looped}{Y}}}$، بهذا الترتيب.
وقد يوقع هذا الاستعمال في التباس؛ فالرموز نفسها نستعملها حين
نتكلم خارج اللغة الصورية على المجموعتين~${\OLId{latin-looped}{X}}$ و${\OLId{latin-looped}{Y}}$.
أما هنا فهي اختصارات لـ!!{formula}s من الرتبة الثانية،
ترد فيها متغيرات العلاقات الأحادية الموضع~${\OLId{latin-looped}{X}}$ و${\OLId{latin-looped}{Y}}$.

\begin{prop}
!!{sentence}~$\lforall[{\OLId{latin-looped}{X}}][\lforall[{\OLId{latin-looped}{Y}}][((\cardle{{\OLId{latin-looped}{X}}}{{\OLId{latin-looped}{Y}}} \land
    \cardle{{\OLId{latin-looped}{Y}}}{{\OLId{latin-looped}{X}}}) \lif \cardeq{{\OLId{latin-looped}{X}}}{{\OLId{latin-looped}{Y}}})]]$ صحيحة منطقيًا.
\end{prop}

\begin{proof}
إرضاء !!{sentence} في !!a{structure}~$\Struct{M}$
يعني أنه لكل مجموعتين جزئيتين~${\OLId{latin-looped}{X}} \subseteq \Domain{{\OLId{latin-looped}{M}}}$
و${\OLId{latin-looped}{Y}} \subseteq \Domain{{\OLId{latin-looped}{M}}}$ يحصل~$\cardeq{{\OLId{latin-looped}{X}}}{{\OLId{latin-looped}{Y}}}$ متى
حصل~$\cardle{{\OLId{latin-looped}{X}}}{{\OLId{latin-looped}{Y}}}$ و$\cardle{{\OLId{latin-looped}{Y}}}{{\OLId{latin-looped}{X}}}$.
وهذا ثابت لـ\emph{أي} مجموعتين~${\OLId{latin-looped}{X}}$ و${\OLId{latin-looped}{Y}}$ بمبرهنة شرودر--برنشتاين.
\end{proof}

\end{document}
